\capitulo{5}{Aspectos relevantes del desarrollo del proyecto}

Este capítulo recoge las decisiones técnicas más relevantes adoptadas durante el desarrollo del sistema, así como los principales retos encontrados y las soluciones implementadas. A diferencia de los capítulos anteriores, centrados en fundamentos y herramientas, este apartado se enfoca en el diseño del sistema y en la justificación de las elecciones realizadas.

Se abordan aspectos clave como la arquitectura del sistema, la elección de formatos de almacenamiento, los mecanismos de captura de eventos y las estrategias utilizadas para garantizar la eficiencia del sistema.

\section{Ciclo de vida del desarrollo}

El desarrollo del proyecto se ha realizado siguiendo un enfoque iterativo e incremental, ampliamente utilizado en el desarrollo de software por su capacidad para adaptarse a cambios y validar soluciones de forma progresiva \cite{pressman2005}.

Inicialmente se llevó a cabo un análisis exploratorio del problema, centrado en comprender las capacidades de la API de Blender y las posibilidades de capturar información del proceso de modelado.

Posteriormente se desarrolló una versión funcional mínima del sistema que permitió validar aspectos críticos como:
\begin{itemize}
    \item La detección de eventos mediante handlers.
    \item El acceso a datos de la escena en tiempo real.
    \item La viabilidad de almacenar los datos generados.
\end{itemize}

A partir de esta base, el sistema fue evolucionando mediante iteraciones sucesivas en las que se incorporaron nuevas métricas, mejoras en el rendimiento y refinamientos en la estructura del código.

\section{Análisis del problema y definición de métricas}

Uno de los principales retos del proyecto consiste en determinar qué información es relevante para analizar el proceso de modelado 3D.

A partir del análisis realizado, se definieron tres categorías principales de datos:

\begin{itemize}
    \item \textbf{Datos temporales}: permiten analizar la evolución del proceso (tiempo total, pausas, fases).
    \item \textbf{Datos espaciales}: describen la interacción en el espacio 3D (posición, velocidad, proximidad).
    \item \textbf{Datos estructurales}: caracterizan la geometría del modelo (vértices, caras, normales).
\end{itemize}

Estas categorías se implementan en el sistema mediante un conjunto de métricas definidas en estructuras como:

\begin{verbatim}
ANALYSIS = {
    "A1": {"label":"Total time","typology":"Temporal"},
    "A4": {"label":"Movement Speed","typology":"Spatial"},
    "A12":{"label":"Model evolution","typology":"Strategy"}
}
\end{verbatim}

Esta clasificación permite estructurar el análisis y facilita la generación de visualizaciones coherentes.

\section{Diseño de la arquitectura del sistema}

El sistema se ha diseñado siguiendo una arquitectura modular, que separa claramente las responsabilidades del sistema en distintos componentes:

\begin{itemize}
    \item \textbf{Captura de eventos}
    \item \textbf{Procesamiento de datos}
    \item \textbf{Almacenamiento}
    \item \textbf{Visualización}
\end{itemize}

Esta separación reduce el acoplamiento y facilita la escalabilidad del sistema, siguiendo principios de diseño de software como la modularidad y la cohesión funcional \cite{gamma1994}.

\section{Decisión de almacenamiento: CSV frente a base de datos}

Una de las decisiones más relevantes del proyecto fue la elección del formato de almacenamiento.

Aunque el uso de bases de datos (SQL o NoSQL) es habitual en sistemas de análisis de datos, en este proyecto se optó por el uso de archivos CSV.

\textbf{Justificación:}
\begin{itemize}
    \item Simplicidad de implementación.
    \item Integración directa con Blender.
    \item Compatibilidad con herramientas de análisis (Python, Excel, etc.).
\end{itemize}

\textbf{Limitaciones:}
\begin{itemize}
    \item Falta de estructura jerárquica.
    \item Menor eficiencia en grandes volúmenes de datos.
\end{itemize}

El uso de CSV es común en procesos de análisis exploratorio de datos debido a su facilidad de uso \cite{data_analysis}.

Un ejemplo de registro generado es:

\begin{verbatim}
USER_ID,TimeStamp,UserX,UserY,UserZ,VertexDelta,ModeChanged
a1b2c3d4,12.34,1.2,0.5,3.1,5,1
\end{verbatim}

\section{Captura de eventos: uso de handlers}

La captura de eventos se implementa mediante los \textit{handlers} de Blender, que permiten ejecutar código cuando se producen cambios en la escena.

Por ejemplo:

\begin{verbatim}
bpy.app.handlers.depsgraph_update_post.append(operator_tracker)
\end{verbatim}

\textbf{Justificación:}
\begin{itemize}
    \item Permiten capturar eventos en tiempo real.
    \item No requieren modificar el flujo de trabajo del usuario.
\end{itemize}

\textbf{Problema:}
Los handlers pueden ejecutarse con alta frecuencia, lo que puede afectar al rendimiento.

\section{Estrategias de eficiencia del sistema}

Uno de los aspectos más críticos del desarrollo ha sido minimizar el impacto del sistema en el rendimiento de Blender.

Para ello se han implementado varias estrategias:

\subsection{Registro basado en cambios (change-based logging)}

En lugar de registrar datos continuamente, el sistema solo genera registros cuando detecta cambios relevantes:

\begin{verbatim}
if not state_changed and not any_operator_flag:
    return None
\end{verbatim}

Esto reduce significativamente el número de operaciones de escritura.

\subsection{Uso de firmas de estado}

Se utiliza una firma del estado de la escena para detectar cambios:

\begin{verbatim}
signature = (verts, ngons, tris, mode, uv_hash)
\end{verbatim}

Solo se registra información cuando esta firma cambia.

\subsection{Control de frecuencia mediante temporizador}

El sistema utiliza un temporizador con intervalo controlado:

\begin{verbatim}
return 0.25
\end{verbatim}

Esto limita la frecuencia de ejecución a 4 veces por segundo.

\subsection{Filtrado de eventos irrelevantes}

Se filtran eventos que no aportan información significativa, evitando ruido en los datos.

\section{Problemas encontrados y soluciones}

Durante el desarrollo se identificaron varios problemas relevantes:

\subsection{Exceso de datos generados}

\textbf{Problema:} el sistema generaba un volumen excesivo de datos.

\textbf{Solución:} uso de logging basado en cambios y eliminación de duplicados.

\subsection{Detección de eventos complejos}

\textbf{Problema:} algunos eventos no se detectan directamente.

\textbf{Solución:} análisis combinado de operadores y estado interno.

\subsection{Impacto en rendimiento}

\textbf{Problema:} ejecución frecuente de handlers.

\textbf{Solución:}
\begin{itemize}
    \item Reducción de frecuencia.
    \item Minimización de cálculos.
    \item Uso de estructuras ligeras.
\end{itemize}

\section{Integración del sistema completo}

El sistema final integra dos addons principales:

\begin{itemize}
    \item \textbf{Data Logger}: encargado de la recopilación de datos.
    \item \textbf{Sistema de análisis}: encargado del procesamiento y visualización.
\end{itemize}

Esta separación permite desacoplar la captura de datos del análisis, facilitando la reutilización del sistema.

\section{Conclusiones del desarrollo}

El desarrollo del sistema ha permitido demostrar la viabilidad de capturar y analizar la interacción del usuario en entornos de modelado 3D.

Las decisiones adoptadas, especialmente en relación con la captura eficiente de eventos y el almacenamiento de datos, han sido fundamentales para garantizar el correcto funcionamiento del sistema.

Asimismo, el proyecto abre nuevas posibilidades en el análisis de procesos creativos en entornos tridimensionales, constituyendo una base sólida para futuras investigaciones.

Una vez descrito el sistema desarrollado y sus funcionalidades, resulta necesario contextualizar la propuesta dentro del panorama actual de investigación. Para ello, en el siguiente capítulo se analizan los principales trabajos relacionados con el estudio del proceso de diseño en entornos tridimensionales.

